\documentclass[11pt,a4paper]{article}
\usepackage[margin=1in]{geometry}
\usepackage{longtable}
\usepackage{hyperref}
\usepackage{titlesec}
\usepackage{parskip}
\titleformat{\section}{\Large\bfseries}{\thesection}{1em}{}
\begin{document}
\begin{titlepage}
\centering
\vspace*{2cm}
{\Huge \textbf{FintechHQ Project Synopsis}\par}
\vspace{0.6cm}
{\Large Professional synopsis for the 2026 project submission cycle\par}
\vspace{1.2cm}
{\large Project Name: FintechHQ\par}
\vspace{0.3cm}
{\large Date: 2026-03-25\par}
\vfill
{\large Repository Basis: \texttt{/Users/abhijith/ResearchProjects/fintechphase2}\par}
\end{titlepage}
\section*{Document Note}
This LaTeX source mirrors the generated documentation for FintechHQ and is provided so the complete report set is preserved in editable source form.
\section{Executive Summary}
FintechHQ is an AI-enabled fintech platform that combines secure user access, financial market data, portfolio visibility, paper trading, algorithmic analysis, and machine-learning services within a modular microservice-oriented architecture. The repository shows a substantial applied engineering effort rather than a narrow prototype: it includes a modern frontend, a transactional backend, a dedicated ML service, persistence models, migrations, containerization, and academic-report-worthy breadth across multiple fintech domains. A particularly important characteristic of the project is that it combines real implementation with staged mock data in some user interfaces, which makes it simultaneously useful as a working demonstration and honest about its present maturity boundaries.
\section{Problem Statement}
Retail and student investors often face fragmented workflows. Market data lives in one tool, technical analysis in another, educational content somewhere else, and security-sensitive account operations in yet another system. FintechHQ addresses this fragmentation by building a unified platform where research, learning, alerts, paper trading, and AI assistance coexist in a single application landscape. The project therefore responds to both a usability gap and an engineering gap: it seeks to show how a fintech platform can be designed with modularity, analytics, and security in mind from the beginning.
\section{Objectives}
The primary objective is to deliver a single, extensible fintech platform for market exploration, secure access, AI-assisted insight generation, and simulated execution. Secondary objectives include demonstrating modern full-stack engineering, separating transactional and ML workloads, supporting future live-broker integration, and providing academic-grade documentation for architecture, modules, and workflows. Another explicit objective visible in the code is support for Indian market assets and a user experience oriented toward investors rather than raw API consumers.
\section{Architecture Overview}
Next.js 16 frontend with React 19, TypeScript, Tailwind CSS, Radix UI primitives, Recharts, Zustand, Framer Motion, and lightweight-charts for interactive market visualization. FastAPI backend that exposes versioned REST APIs for authentication, user management, stocks, strategies, trading, alerts, community, and education workflows. Dedicated FastAPI ML microservice that handles feature engineering, price prediction, recommendation, risk analytics, fundamental analysis, backtesting, sentiment scoring, and AI chart-pattern detection. PostgreSQL with TimescaleDB-style time-series usage for market and OHLCV storage, supported by Alembic migrations and SQLAlchemy ORM models. Redis for caching, fast state lookups, and service-side infrastructure support. Docker Compose orchestration for frontend, backend, ML service, database, Redis, MLflow, and Flower.
\section{Core Strengths}
Strong authentication and security posture with Argon2 password hashing, password-complexity enforcement, breached-password checks, JWT access and refresh tokens, lockout handling, password history, and optional TOTP 2FA. Broad functional coverage across onboarding, portfolio views, market data, technical analysis, fundamentals, news, community, strategies, backtesting, alerts, learning, and AI-assisted workflows. Clear separation between transactional business logic and compute-heavy analytics through a dedicated ML service. Indian market orientation through curated Nifty and Sensex stock universes, Indian ticker handling, and rupee-oriented presentation in several UI modules. Real-time capability through WebSocket channels for simulated prices, user-specific portfolio streams, and alert notifications.
\section{Current Implementation Status}
The codebase contains a mix of production-style backend services and UI screens that are either API-integrated or still backed by representative mock data. Authentication, password reset, stock search, quotes, OHLCV retrieval, indicator calculation, alerts, community posting, education seeding, and several ML endpoints are implemented as concrete service flows. Some screens such as portfolio, trade, risk, and backtest currently render realistic demo data in the frontend while the backend and ML layers already expose foundations for fuller integration. The project uses legacy naming in places such as 'FinTech AI Platform' and older repository documents; this document set standardizes the academic/project title as FintechHQ, as requested.
\section{Authentication and Identity}
The authentication and identity module addresses secure onboarding, login, session continuity, password recovery, and account hardening. Within FintechHQ, this area is not treated as a standalone academic idea; it is directly represented in the codebase through dedicated UI routes, backend endpoints, services, models, and data contracts. That makes the module suitable for both demonstration and further extension. The current implementation shows a well-developed backend flow with registration, login, refresh tokens, account lockouts, breached-password checks, password history, and optional TOTP-based 2FA, along with polished login, register, forgot-password, reset-password, and 2FA setup screens. The project architecture deliberately separates presentation concerns from domain logic, so user actions initiated in the frontend can be routed through the backend for validation and persistence while the ML service handles heavier analytics where needed. From a systems perspective, credentials enter through the auth UI, are submitted to the FastAPI auth routes, pass password-policy and compromise checks, are hashed with Argon2, and then return token material and session state for protected application access. This layered flow is one of the strongest structural qualities of the project because it supports future replacement of mock screens with live data without forcing a redesign of the entire product. Operationally and from a security standpoint, the project enforces rate limits, stores session metadata, tracks failed attempts, and blocks insecure password reuse during resets. This is important in a fintech setting because user trust, auditability, and predictable behavior are as important as raw feature count. Looking ahead, production readiness would improve further by aligning the frontend 2FA setup page with the real backend endpoints and moving token storage toward httpOnly-cookie strategies if the team adopts that model. The existing code already contains enough hooks and service boundaries to support that growth path with incremental engineering work rather than a full rewrite.
\section{Dashboard and Portfolio Intelligence}
The dashboard and portfolio intelligence module addresses executive visibility into account value, portfolio allocation, alerts, market movers, and AI-oriented summaries. Within FintechHQ, this area is not treated as a standalone academic idea; it is directly represented in the codebase through dedicated UI routes, backend endpoints, services, models, and data contracts. That makes the module suitable for both demonstration and further extension. The current implementation shows a composed dashboard built from reusable React components such as balance cards, sparkline visualizations, allocation charts, alert banners, and market-mover widgets, plus a portfolio page that demonstrates holdings, allocation, and recent activity flows. The project architecture deliberately separates presentation concerns from domain logic, so user actions initiated in the frontend can be routed through the backend for validation and persistence while the ML service handles heavier analytics where needed. From a systems perspective, data can be assembled from portfolio state, market data, alerts, and enrichment services, then surfaced as summary cards and drill-down views for faster user decision-making. This layered flow is one of the strongest structural qualities of the project because it supports future replacement of mock screens with live data without forcing a redesign of the entire product. Operationally and from a security standpoint, the architecture supports role-based protection, authenticated data fetching, and clear separation between market data retrieval and user-specific financial state. This is important in a fintech setting because user trust, auditability, and predictable behavior are as important as raw feature count. Looking ahead, the remaining work is mainly integration depth: replacing representative mock portfolio values with live backend portfolio summaries and extending analytics over real transaction history. The existing code already contains enough hooks and service boundaries to support that growth path with incremental engineering work rather than a full rewrite.
\section{Market Data and Technical Analysis}
The market data and technical analysis module addresses price discovery, indicator analysis, chart interpretation, and symbol-level exploration. Within FintechHQ, this area is not treated as a standalone academic idea; it is directly represented in the codebase through dedicated UI routes, backend endpoints, services, models, and data contracts. That makes the module suitable for both demonstration and further extension. The current implementation shows real stock search, quote retrieval, OHLCV fetching, server-side indicator generation, ticker autocomplete, charting through lightweight-charts, and an AI pattern-detection panel connected to the ML stack. The project architecture deliberately separates presentation concerns from domain logic, so user actions initiated in the frontend can be routed through the backend for validation and persistence while the ML service handles heavier analytics where needed. From a systems perspective, the user selects a symbol and timeframe, the frontend calls protected stock endpoints, the backend serves candle and quote data from stored or live sources, and the ML service can inspect recent candles for advanced patterns. This layered flow is one of the strongest structural qualities of the project because it supports future replacement of mock screens with live data without forcing a redesign of the entire product. Operationally and from a security standpoint, route ordering, authenticated requests, and modular service boundaries reduce errors while keeping extension points open for richer indicators and broker-grade data feeds. This is important in a fintech setting because user trust, auditability, and predictable behavior are as important as raw feature count. Looking ahead, future growth can include broader timeframe coverage, deeper indicator libraries, exchange-grade feeds, and stronger synchronization between chart overlays and pattern annotations. The existing code already contains enough hooks and service boundaries to support that growth path with incremental engineering work rather than a full rewrite.
\section{Fundamental, News, and Research Intelligence}
The research intelligence module addresses fundamental valuation, company insight, news monitoring, and sentiment-enriched decision support. Within FintechHQ, this area is not treated as a standalone academic idea; it is directly represented in the codebase through dedicated UI routes, backend endpoints, services, models, and data contracts. That makes the module suitable for both demonstration and further extension. The current implementation shows a fundamental analysis page that displays valuation, health score, margins of safety, ratios, and assumptions, plus a news page that fetches asset-level stories and shows sentiment badges. The project architecture deliberately separates presentation concerns from domain logic, so user actions initiated in the frontend can be routed through the backend for validation and persistence while the ML service handles heavier analytics where needed. From a systems perspective, stock-level research requests move through backend routes into market-data and enrichment services, after which the results are displayed in structured cards and analyst-style summaries. This layered flow is one of the strongest structural qualities of the project because it supports future replacement of mock screens with live data without forcing a redesign of the entire product. Operationally and from a security standpoint, the project benefits from typed API responses, authenticated access, and clear demarcation between raw provider data, transformed metrics, and presentational UI cards. This is important in a fintech setting because user trust, auditability, and predictable behavior are as important as raw feature count. Looking ahead, future versions can strengthen this area by persisting research snapshots, adding analyst consensus history, and improving narrative generation for company-level reports. The existing code already contains enough hooks and service boundaries to support that growth path with incremental engineering work rather than a full rewrite.
\section{Trading, Strategies, Alerts, and Community}
The active-operations module addresses paper trading, strategy definition, alerting, and investor interaction. Within FintechHQ, this area is not treated as a standalone academic idea; it is directly represented in the codebase through dedicated UI routes, backend endpoints, services, models, and data contracts. That makes the module suitable for both demonstration and further extension. The current implementation shows a broker simulation service that creates portfolios, validates buying power, updates positions, records transactions, and handles order states, alongside strategy CRUD, alert management, notifications, and social posting/commenting flows. The project architecture deliberately separates presentation concerns from domain logic, so user actions initiated in the frontend can be routed through the backend for validation and persistence while the ML service handles heavier analytics where needed. From a systems perspective, orders originate in the frontend, flow into trading endpoints, are validated against current simulated prices and cash, and then update portfolio records and transaction logs while alerts and social flows run in parallel through their own services. This layered flow is one of the strongest structural qualities of the project because it supports future replacement of mock screens with live data without forcing a redesign of the entire product. Operationally and from a security standpoint, auditability is supported through transaction records, order status tracking, notification history, and protected user-specific endpoint dependencies. This is important in a fintech setting because user trust, auditability, and predictable behavior are as important as raw feature count. Looking ahead, additional realism can be achieved with live market feeds, richer order types, event-driven order execution, and deeper community moderation and ranking logic. The existing code already contains enough hooks and service boundaries to support that growth path with incremental engineering work rather than a full rewrite.
\section{Machine Learning, Risk, and Backtesting}
The AI and quantitative analytics module addresses feature generation, forecasting, recommendations, risk measurement, chart-pattern recognition, sentiment scoring, and simulated strategy evaluation. Within FintechHQ, this area is not treated as a standalone academic idea; it is directly represented in the codebase through dedicated UI routes, backend endpoints, services, models, and data contracts. That makes the module suitable for both demonstration and further extension. The current implementation shows a dedicated ML microservice with separate endpoints and services for prediction, recommendation, risk, backtesting, sentiment, feature engineering, fundamental analytics, training, and classical plus deep-learning pattern detection. The project architecture deliberately separates presentation concerns from domain logic, so user actions initiated in the frontend can be routed through the backend for validation and persistence while the ML service handles heavier analytics where needed. From a systems perspective, the backend or frontend can submit structured market data to the ML service, where numeric preprocessing, inference, calibration, and optional annotated visualization produce analytics outputs for the user interface. This layered flow is one of the strongest structural qualities of the project because it supports future replacement of mock screens with live data without forcing a redesign of the entire product. Operationally and from a security standpoint, the service design keeps compute-heavy logic away from the transactional API and supports traceability through explicit endpoint boundaries and model-oriented service classes. This is important in a fintech setting because user trust, auditability, and predictable behavior are as important as raw feature count. Looking ahead, the next maturity stage would include stronger model lifecycle management, repeatable training pipelines, stored experiment metadata, richer validation datasets, and more complete live integration into currently mocked UI screens. The existing code already contains enough hooks and service boundaries to support that growth path with incremental engineering work rather than a full rewrite.
\section{Technology Stack}
Frontend uses Next.js 16 app router, React 19 and TypeScript, Tailwind CSS 4, Radix UI and shadcn-style components, Recharts and lightweight-charts, NextAuth scaffold and local token storage flow, Framer Motion and tsparticles for interactive UI polish. Backend uses FastAPI, SQLAlchemy 2 async ORM, Alembic migrations, PostgreSQL / TimescaleDB usage pattern, Redis, python-jose for JWT handling, passlib with Argon2, sha256\_crypt, and bcrypt support, slowapi for rate limiting, yfinance, pandas, numpy, ta. ML Service uses FastAPI ML microservice, PyTorch, scikit-learn, transformers, MLflow, OpenCV headless, Pillow, mplfinance, SciPy and custom algorithmic signal processing. DevOps uses Docker Compose, Dockerfiles per service, MLflow tracking container, Flower monitoring container, Container networking and mounted development volumes.
\section{Security Architecture}
Security is one of the clearest strengths of the repository. The backend not only hashes passwords and issues tokens, but also applies complexity rules, breached-password screening, session recording, lockout behavior, password history enforcement, and optional TOTP verification. For a student or portfolio project, that breadth meaningfully increases the credibility of the platform because it shows awareness of real fintech expectations.
\section{Data Model and Persistence}
The backend data model spans identity, sessions, portfolio accounting, order management, alerts, community, learning content, market data, ML output tracking, watchlists, and news enrichment. This breadth enables the platform to behave like a connected product instead of a thin dashboard over external APIs. The model layer also reflects sensible domain separation, which will support migration, reporting, and future analytics work.
\section{Functional Requirements}
Users shall be able to register, log in, refresh sessions, reset passwords, and optionally enable 2FA. Users shall be able to search stocks and access symbol-level metadata. Users shall be able to view historical OHLCV data and server-generated indicators. Users shall be able to request quotes, market movers, fundamentals, and news. Users shall be able to place paper-trading orders and inspect positions. Users shall be able to create and manage alerts and receive notifications. Users shall be able to create and manage strategies. Users shall be able to access educational modules and track lesson progress. Users shall be able to create community posts and comments. The platform shall expose ML endpoints for prediction, risk, recommendation, sentiment, backtesting, and pattern detection. The platform shall provide live or near-live interactive experiences through WebSocket subscriptions where applicable. The system shall support Indian-market symbol discovery through curated exchange lists.
\section{Non-Functional Requirements}
Security controls shall be strong enough for a fintech demonstration environment. The architecture shall remain modular and separable by service boundary. The system shall be locally reproducible through containerized deployment. The codebase shall support API versioning and maintainable router organization. The data model shall support growth across users, portfolios, alerts, strategies, and learning content. The UI shall remain understandable across multiple financial workflows. Analytics workloads shall be isolated from the transactional backend. The system shall be extensible for future live-broker and live-feed integration. The project shall remain demonstrable even when some views use mock data. The documentation shall distinguish clearly between implemented and staged functionality.
\section{Backend API Design}
The REST surface is organized by domain, which makes the codebase readable and maintainable. Authentication, users, stocks, trading, alerts, strategies, social, and education are each isolated behind their own routers and services. This pattern reduces coupling and makes testing and future refactoring easier.
\section{ML Service Design}
The ML service is not just a placeholder route collection; it contains distinct service classes for prediction, recommendation, risk, feature engineering, fundamental analysis, training, sentiment, backtesting, chart rendering, and pattern detection. The pattern-detection subsystem is especially notable because it combines algorithmic detection with a CNN-plus-Transformer pathway and keeps recent history for inspection. This gives the platform a strong AI narrative backed by code, not only by documentation.
\section{Real-Time and Interactive Features}
The backend exposes a WebSocket endpoint with subscription semantics for price, portfolio, and alert channels. This creates a route toward real-time experiences even where current data streams are simulated. In platform terms, the WebSocket layer is important because it shows how the application can evolve from periodic refresh to event-driven fintech interaction.
\section{Deployment Topology}
frontend container: Next.js development server exposed on port 3000 with API and ML-service environment variables. backend container: FastAPI core API exposed on port 8000 with database and Redis connectivity. ml-service container: FastAPI ML analytics service exposed on port 8001 with model-volume persistence. db container: PostgreSQL / TimescaleDB-oriented storage service on port 5432. redis container: In-memory cache and broker service on port 6379. mlflow container: Experiment tracking service exposed through port 5001. flower container: Task-queue monitoring interface exposed on port 5555.
\section{Testing Perspective and Quality Gaps}
The repository includes structural implementation depth but does not expose a strong automated test suite in the inspected files. Manual verification paths appear to be the current dominant validation mode. The presence of service separation improves testability even where explicit tests are not yet extensive. Security-sensitive flows deserve integration tests around token issuance, reset behavior, and 2FA completion. Trading and alert workflows would benefit from deterministic simulation tests. ML endpoints would benefit from fixture-driven inference and contract tests. Frontend modules that still use mock data should be validated again after API integration. A mature next step would add CI checks for linting, backend tests, and typed frontend builds.
\section{Use Cases}
User onboarding: A user registers with strong credentials, optionally configures 2FA, receives tokens, and enters the dashboard. Market research: The user searches Indian symbols, loads OHLCV data, studies indicators, checks quotes, reads news, and explores fundamentals. Technical analysis: The technical-analysis screen fetches candles, indicator overlays, quotes, and ML-backed chart-pattern findings. Paper trading: Trading logic validates buying power, updates positions, records transactions, and maintains order state. Portfolio monitoring: The portfolio summary and dashboard views present holdings, P\&L, allocation, alerts, and news context. Alerting: Users configure price thresholds that can later produce notifications and alert-history entries. Backtesting and strategies: Users design strategies, run simulations, and store strategy definitions for iterative refinement. AI-assisted analytics: Risk, prediction, recommendation, sentiment, and pattern-detection endpoints enrich the investment workflow.
\section{Limitations and Engineering Risks}
Several UI pages still render mock or staged data even though backend service foundations are present. Reset-password flow exposes a reset token in the response for demo convenience and should not do so in production. The 2FA setup page currently uses mocked frontend data rather than invoking the real backend setup endpoint. The backend configuration defines an RS256 label, while token creation and decode helpers currently use HS256 symmetric signing internally. Real broker integration is not implemented; trading is a paper-trading simulation. Pattern image detection depends on trained YOLO or related weights and currently documents that limitation in the API response.
\section{Future Scope}
The project is well positioned for live broker APIs, stronger ML lifecycle management, richer database seeding, real-time exchange feeds, improved secrets management, and fuller frontend-backend wiring for every screen. Additional work can also focus on production observability, automated test coverage, authorization policies, and deployment hardening for cloud environments. Because FintechHQ already has clear service boundaries, this future scope is incremental and realistic.
\section{Conclusion}
FintechHQ is a strong full-stack fintech project with meaningful breadth and a credible architecture. It already demonstrates secure access, market-data workflows, AI-driven analytics, paper-trading logic, and modular growth potential. For academic evaluation, it provides enough substance to support a synopsis, multi-stage reports, and multiple presentation tiers without relying on invented capabilities.
\section*{Appendix A: Security Controls}
\begin{itemize}
\item Access and refresh JWT tokens with explicit token type claims.
\item Argon2 password hashing with compatibility fallbacks for legacy hash migration.
\item Password policy requiring length, mixed case, numerics, and special characters.
\item Breached-password screening using HaveIBeenPwned k-anonymity API logic.
\item Failed-login tracking and temporary account lockout after repeated failures.
\item User session recording with device fingerprint and IP address capture.
\item Password history enforcement during password reset.
\item TOTP secret generation and code verification for two-factor authentication.
\item SlowAPI-based rate limiting on registration, login, and 2FA endpoints.
\item CORS controls and authenticated dependency guards for protected APIs.
\end{itemize}
\section*{Appendix B: Data Models}
\begin{longtable}{p{0.28\textwidth} p{0.64\textwidth}}
\textbf{Item} & \textbf{Description} \\ \hline
User & Identity, account status, risk profile, verification state, lockouts, timestamps, and security relationships. \\
PasswordHistory & Tracks previous password hashes to block reuse. \\
UserSession & Stores per-device session tokens, expiry, fingerprint, and IP metadata. \\
Stock / OHLCVData / FundamentalData & Market master data, historical candles, and stored fundamentals. \\
Portfolio / Position / Transaction & Cash balance, holdings, and transaction ledger for paper trading. \\
Order / Trade & Orders, execution status, fills, and trading activity. \\
Strategy / BacktestRun & User-defined strategies and simulation results. \\
Alert / Notification / AlertHistory & Threshold definitions, user alerts, read-state tracking, and trigger logs. \\
Post / Comment / Like / Follow & Community interaction graph and social feed data. \\
Module / Lesson / UserProgress / Enrollment / Quiz / QuizSubmission & Learning management and progress tracking. \\
Watchlist / WatchlistItem & User-defined monitoring lists. \\
MLPrediction / NewsArticle / NewsSentiment & Predictive outputs and enriched news records. \\
\end{longtable}
\section*{Appendix C: Backend API Inventory}
\begin{longtable}{p{0.28\textwidth} p{0.64\textwidth}}
\textbf{Item} & \textbf{Description} \\ \hline
/api/v1/auth/register & Registers a user after uniqueness checks, password policy validation, breach screening, and password-history initialization. \\
/api/v1/auth/login/access-token & Authenticates users, enforces lockouts, supports 2FA branching, and issues access plus refresh tokens. \\
/api/v1/auth/login/verify-2fa & Completes a TOTP-based login challenge and opens a real session. \\
/api/v1/auth/refresh & Refreshes access tokens using refresh-token semantics. \\
/api/v1/auth/2fa/setup & Generates a TOTP secret and provisioning URI. \\
/api/v1/auth/2fa/verify & Finalizes 2FA setup. \\
/api/v1/users/me & Returns the authenticated user profile. \\
/api/v1/users/forgot-password & Generates a reset token with demo-mode response behavior. \\
/api/v1/users/reset-password & Validates reset token, enforces password rules, blocks reuse, and invalidates sessions. \\
/api/v1/stocks/search & Searches stored stocks and supplements results from the curated Indian stock universe. \\
/api/v1/stocks/market-movers & Returns gainers and losers using live or fallback data. \\
/api/v1/stocks/index-stocks & Supplies Nifty, Sensex, or combined symbol catalogs. \\
/api/v1/stocks/\{ticker\} & Returns or hydrates stock master information. \\
/api/v1/stocks/\{ticker\}/ohlcv & Provides candle data from DB or live provider fallback. \\
/api/v1/stocks/\{ticker\}/indicators & Calculates SMA, EMA, RSI, MACD, and Bollinger Bands server-side. \\
/api/v1/stocks/\{ticker\}/quote & Supplies latest quote data. \\
/api/v1/stocks/\{ticker\}/fundamentals & Returns fundamentals and valuation-oriented data. \\
/api/v1/stocks/\{ticker\}/news & Fetches and sentiment-enriches asset-level news. \\
/api/v1/trading/orders & Places or lists paper-trading orders. \\
/api/v1/trading/orders/\{order\_id\} & Cancels pending paper orders. \\
/api/v1/trading/positions & Returns holdings enriched with simulated live prices. \\
/api/v1/trading/portfolio/summary & Provides aggregate portfolio values. \\
/api/v1/alerts & Creates, lists, and deletes alerts. \\
/api/v1/alerts/notifications & Lists and marks notifications as read. \\
/api/v1/strategies & Lists, creates, and deletes strategy definitions. \\
/api/v1/social/feed & Returns social feed content. \\
/api/v1/social/posts & Creates community posts. \\
/api/v1/social/posts/\{id\}/comments & Adds comments to community posts. \\
/api/v1/education/seed & Seeds learning content. \\
/api/v1/education/modules & Lists modules and learning state. \\
/api/v1/education/lessons/\{id\}/complete & Marks lessons complete. \\
\end{longtable}
\section*{Appendix D: ML API Inventory}
\begin{longtable}{p{0.28\textwidth} p{0.64\textwidth}}
\textbf{Item} & \textbf{Description} \\ \hline
/api/v1/features/calculate & Generates technical features, candlestick patterns, and support/resistance levels. \\
/api/v1/training/train & Runs model-training related workflows. \\
/api/v1/sentiment/predict & Scores text sentiment for financial context. \\
/api/v1/prediction/train & Trains an LSTM-based price model. \\
/api/v1/prediction/forecast & Forecasts the next price from recent history. \\
/api/v1/recommendation/analyze & Produces a recommendation from price behavior and optional news text. \\
/api/v1/technical/analyze & Runs technical-analysis logic for chart-based insight. \\
/api/v1/risk/analyze & Calculates portfolio risk metrics from holdings and history. \\
/api/v1/fundamental/analyze & Computes DCF and health-score outputs. \\
/api/v1/backtest/run & Executes a simplified backtesting routine. \\
/api/v1/pattern-detection/detect-from-data & Runs ensemble pattern detection and optionally returns annotated charts. \\
/api/v1/pattern-detection/detect-image & Placeholder image-based path pending trained vision weights. \\
/api/v1/pattern-detection/patterns-history & Returns stored detection history. \\
/api/v1/pattern-detection/supported-patterns & Lists supported classical chart patterns. \\
\end{longtable}
\section*{Appendix E: Frontend Route Inventory}
\begin{longtable}{p{0.28\textwidth} p{0.64\textwidth}}
\textbf{Item} & \textbf{Description} \\ \hline
/login & Secure login form with password visibility toggle, API authentication, and 2FA branching. \\
/register & Three-step registration flow capturing identity, password quality, and investment profile. \\
/forgot-password & Password-reset request screen that triggers reset-token generation. \\
/reset-password & Protected reset flow with password-strength enforcement. \\
/2fa-setup & Two-factor setup UI with QR-code style onboarding and code verification pattern. \\
/dashboard & Executive landing page with portfolio balance, allocation, alerts, movers, and AI briefing widgets. \\
/portfolio & Holdings, allocation, and recent portfolio activity view. \\
/technical & Interactive charting, quotes, indicators, and AI chart-pattern interface. \\
/fundamental & Valuation, health score, and ratio-based company research screen. \\
/news & Asset-specific news feed with sentiment badges. \\
/trade & Paper-trading order ticket, order book, order history, and position views. \\
/strategies & Strategy list and strategy-creation entry points. \\
/backtest & Simulation workspace for strategy evaluation. \\
/risk & Portfolio risk metrics and AI-generated interpretation panel. \\
/alerts & Alert creation and alert management screen. \\
/community & Social or community discussion and posting area. \\
/learn & Learning-center interface for finance education content. \\
\end{longtable}
\section*{Appendix F: Backend Services}
\begin{longtable}{p{0.28\textwidth} p{0.64\textwidth}}
\textbf{Item} & \textbf{Description} \\ \hline
StockDataService & Fetches security metadata, historical data, live quotes, market movers, fundamentals, index data, and enriched news. \\
BrokerService & Creates default portfolios, validates orders, simulates fills, updates positions, tracks cash, and handles order cancellation. \\
AlertService & Creates alerts, serves notifications, tracks read state, and triggers alert history when price conditions are met. \\
EducationService & Seeds course content, lists modules, and tracks lesson completion. \\
SocialService & Creates posts and comments and retrieves social feed data. \\
WebSocket ConnectionManager & Maintains live client connections, subscriptions, broadcasts, and personal messages. \\
\end{longtable}
\section*{Appendix G: ML Services}
\begin{longtable}{p{0.28\textwidth} p{0.64\textwidth}}
\textbf{Item} & \textbf{Description} \\ \hline
FeatureService & Calculates technical indicators, candlestick patterns, and support/resistance levels. \\
PredictionService & Trains and serves an LSTM-based next-price forecaster. \\
RecommendationService & Builds recommendation logic from raw data and optional news context. \\
RiskService & Computes VaR, volatility, Sharpe ratio, and maximum drawdown from aligned return series. \\
FundamentalService & Calculates DCF outputs and company health scores. \\
BacktestService & Runs simplified backtesting workflows for strategy evaluation. \\
SentimentService & Loads a transformer-style sentiment model and predicts financial text sentiment. \\
TrainingService & Coordinates training-oriented workflows. \\
PatternDetectionService & Combines algorithmic and deep-learning chart-pattern detection. \\
ChartRenderer & Renders and annotates chart images for pattern output. \\
PatternHistoryStore & Stores recent pattern-detection history for later retrieval. \\
\end{longtable}
\section*{Appendix H: Container Topology}
\begin{longtable}{p{0.28\textwidth} p{0.64\textwidth}}
\textbf{Item} & \textbf{Description} \\ \hline
frontend & Next.js development server exposed on port 3000 with API and ML-service environment variables. \\
backend & FastAPI core API exposed on port 8000 with database and Redis connectivity. \\
ml-service & FastAPI ML analytics service exposed on port 8001 with model-volume persistence. \\
db & PostgreSQL / TimescaleDB-oriented storage service on port 5432. \\
redis & In-memory cache and broker service on port 6379. \\
mlflow & Experiment tracking service exposed through port 5001. \\
flower & Task-queue monitoring interface exposed on port 5555. \\
\end{longtable}
\end{document}
